\documentclass[twocolumn]{article}

\usepackage{graphicx}
\usepackage{booktabs}
\usepackage{geometry}
\usepackage{hyperref}
\usepackage{minted}
\usepackage{xcolor}
\usepackage{float}
\usepackage{amsmath}

\geometry{a4paper, margin=0.75in}

\hypersetup{
    colorlinks=true,
    linkcolor=blue,
    filecolor=magenta,
    urlcolor=blue,
    citecolor=blue,
}

\definecolor{bg}{rgb}{0.95,0.95,0.95}

\title{\textbf{Pothole Detection and Classification: A Data-Centric Approach using Improved YOLOv8}}
\date{February 2026}

\begin{document}

\maketitle

\begin{abstract}
\noindent Real-world deployment of object detection models for infrastructure monitoring is often hampered by a high rate of false positives. This study presents a data-centric methodology to enhance the robustness of a two-stage YOLOv8-based pothole detection and severity classification system. The baseline detection model, while capable of identifying defects, exhibited poor specificity. The intervention involved retraining this model on a refined dataset augmented with negative samples (images of undamaged roads). This paper details the experimental setup, model parameters, and a comparative analysis of the results. The refined detection model achieved a precision of 94.7\%, a dramatic improvement that validates the efficacy of negative sample training. We also detail the secondary classification model used for severity analysis and its performance.
\end{abstract}

\section{Introduction}
The maintenance of road infrastructure is a critical and costly task for municipalities worldwide. Potholes and other surface defects not only degrade driver comfort but also pose significant safety risks. Traditional methods of road inspection rely on manual surveys, which are labor-intensive and slow. Computer vision offers a scalable alternative, and the YOLOv8 architecture \cite{yolov8_ultralytics} is a state-of-the-art detector well-suited for this task. However, a primary challenge is ensuring model reliability. Models trained without a sufficiently diverse dataset often generate a high rate of false positives. This study details a successful data-centric approach to overcome this challenge.

\section{Experimental Analysis}

\subsection{Experimental Setup}
All training procedures were executed on a workstation with the following configuration:
\begin{itemize}
    \item \textbf{GPU:} NVIDIA RTX 3050
    \item \textbf{Frameworks:} PyTorch 2.6.0 \cite{NEURIPS2019_9015}
    \item \textbf{Primary Models:} YOLOv8m (Detection), YOLOv8x-cls (Classification)
\end{itemize}

\subsection{Model Parameters and Data}
The core intervention was a strategic shift in data curation. Table \ref{tab:intervention} compares the key parameters of the baseline and refined detection models.

\begin{table}[H]
\centering
\caption{Comparison of Baseline and Refined Model Configurations}
\label{tab:intervention}
\resizebox{\columnwidth}{!}{%
\begin{tabular}{@{}lll@{}}
\toprule
\textbf{Parameter} & \textbf{Baseline System} & \textbf{Refined System} \\ \midrule
\textbf{Architecture} & YOLOv8x (Larger) & YOLOv8m (Smaller) \\
\textbf{Training Data} & Positive examples only & \begin{tabular}[c]{@{}l@{}}Positive examples +\\ \textbf{651 negative samples}\end{tabular} \\
\textbf{Task Focus} & Multi-class detection & Single-class (pothole) \\ \bottomrule
\end{tabular}%
}
\end{table}

\section{Analysis of False Positive Reduction: A Before-and-After Comparison}
This section details the behavior of the detection model before and after the data-centric intervention, explaining the cause of the initial problem and the reason for the subsequent improvement.

\subsection{Phase 1: The Baseline Model's Behavior}
The initial model (`YOLOv8x`) was trained on a dataset composed exclusively of images that contained potholes.

\subsubsection{Behavior and Cause of False Positives}
Because the model never saw examples of what it was supposed to *ignore*, its understanding of a "pothole" was overly simplistic. It learned to associate generic features like "a dark patch on a grey surface" or "an area of irregular texture" with the pothole class.

This led to highly problematic behavior in practice. When the model analyzed an image of a perfectly healthy road that had a simple shadow, a water stain, a dark-colored repair patch, or even just textured asphalt, it would incorrectly classify these benign features as potholes. This high rate of false positives made the model unreliable for any automated process, as it created too much "noise" for manual verification. While specific accuracy metrics for this phase are unavailable from the project logs, its qualitative failure was the documented catalyst for this study.

\subsection{Phase 2: The Refined Model's Behavior}
The intervention was simple: we retrained the model on a new dataset that included **651 images with no potholes**. These "negative samples" were the key to changing the model's behavior.

\subsubsection{Behavioral Change and Reduction of False Positives}
By adding negative samples, we forced the model to learn a more complex and nuanced task. It now had to learn not only what a pothole *is*, but just as importantly, what a pothole *is not*.

During the new training process, every time the model would have incorrectly predicted a pothole on a clean road image, the training algorithm would penalize it. This continuous correction forced the model's internal feature extractors to become more discerning. It could no longer rely on simple cues like "dark patch." It had to learn the more specific, combined features of a true pothole (e.g., sharp edges, specific shadow patterns indicating depth, and surrounding surface deformation). The model learned to suppress the feature patterns associated with healthy roads, effectively ignoring them.

\subsubsection{The Outcome: A Dramatic Increase in Accuracy}
This change in behavior is directly and quantitatively measured by the final accuracy metrics of the refined model, shown in Table \ref{tab:performance_refined}.

\begin{table}[H]
\centering
\caption{Validation Performance of the Refined M1 Detection Model}
\label{tab:performance_refined}
\begin{tabular}{@{}lc@{}}
\toprule
\textbf{Metric} & \textbf{Value} \\ \midrule
\textbf{Precision} & \textbf{94.7\%} \\
Recall & 91.6\% \\
mAP@.50 & 95.7\% \\
mAP@.50-.95 & 81.5\% \\ \bottomrule
\end{tabular}
\end{table}

The most critical result is the **Precision of 94.7\%**. Precision is the metric that tracks false positives ($\text{TP} / (\text{TP} + \text{FP})$). This high score is the direct result of the model learning to ignore non-pothole features. This single change transformed the system from an unreliable prototype into a highly accurate and deployable tool.

\section{System Architecture}
The final system is a two-stage pipeline designed for both accuracy and actionable insights. This modular design allows each model to be highly specialized for its task.

\begin{enumerate}
    \item \textbf{M1 - Detection Model:} The refined YOLOv8m model analyzes the full input image to identify and locate all potential potholes.
    \item \textbf{M2 - Classification Model:} For each detected pothole, the image region is cropped and passed to a YOLOv8x-cls model, which categorizes its severity.
\end{enumerate}

\subsection{Severity Classification Model (M2)}
To provide practical utility, the M2 model was trained to classify detected potholes into three levels of severity.

\subsubsection{Defining Severity Classes}
The dataset for the M2 model was labeled based on the following criteria:
\begin{itemize}
    \item \textbf{Minor:} Surface-level defects and shallow depressions unlikely to cause immediate vehicle damage.
    \item \textbf{Medium:} Formed potholes with visible depth that pose a risk of tire or suspension wear.
    \item \textbf{Major:} Large and/or deep potholes that represent a significant and immediate safety hazard.
\end{itemize}

\subsubsection{M2 Model Performance}
The M2 model is a YOLOv8x-cls architecture trained on pre-cropped images of potholes. It learns to differentiate severity by analyzing features like texture and shadow that indicate depth. Its performance is summarized by the confusion matrix in Figure \ref{fig:sev_conf_matrix}, which shows a strong ability to distinguish between the classes.

\begin{figure}[H]
    \centering
    \includegraphics[width=\columnwidth]{confusion_matrix_normalized.png}
    \caption{Normalized Confusion Matrix for the M2 Severity Classification Model. The strong diagonal indicates good performance in distinguishing between the three severity classes.}
    \label{fig:sev_conf_matrix}
\end{figure}

\subsection{Dataset Description}

For training and evaluation of the pothole detection models, multiple publicly available datasets were used to ensure diversity in road conditions, lighting variations, and annotation quality. The datasets were partitioned into training and validation sets using an 80/20 split, ensuring that the models were evaluated on a hold-out set of unseen data.

\begin{itemize}

\item \textbf{General Pothole Dataset (Used for Model M1)}  
The primary dataset used for Model M1 was sourced from Kaggle and includes a wide variety of pothole images captured under real-world conditions.  
\newline
Dataset Link: \url{https://www.kaggle.com/datasets/srajanchourasia/pothole-dataset}

\item \textbf{Pothole Detection Dataset Repository (Used for Model M1)}  
Additional reference data and structured dataset information were obtained from Dataset Ninja to improve data understanding and preprocessing consistency.  
\newline
Dataset Link: \url{https://datasetninja.com/pothole-detection}

\item \textbf{Annotated Potholes with Severity Levels (Used for Model M2)}  
For Model M2, a severity-annotated pothole dataset was used. This dataset contains pothole images labeled according to severity levels, enabling the model to learn not only detection but also pothole condition classification.  
\newline
Dataset Link: \url{https://www.kaggle.com/datasets/idanbaru/annotated-potholes-with-severity-levels}

\end{itemize}

These datasets collectively improved model robustness by introducing variation in road texture, environmental conditions, and annotation granularity.

\subsection{Dataset Composition}
The refined dataset for the M1 pothole detection model was constructed by combining positive samples (images containing potholes) with negative samples (images of undamaged roads). This strategy is crucial for training a model that can accurately distinguish between actual road defects and benign features like shadows or stains. The detailed composition of this dataset is shown in Table \ref{tab:dataset_composition}.

\begin{table}[H]
\centering
\caption{Composition of the Refined Pothole Detection Dataset (M1)}
\label{tab:dataset_composition}
\resizebox{\columnwidth}{!}{%
\begin{tabular}{@{}lccc@{}}
\toprule
\textbf{Data Type} & \textbf{Training Set} & \textbf{Validation Set} & \textbf{Total} \\ \midrule
Pothole Images (Positive) & 1,097 & 285 & 1,382 \\
Good Road Images (Negative) & 520 & 131 & 651 \\ \midrule
\textbf{Total} & \textbf{1,617} & \textbf{416} & \textbf{2,033} \\ \bottomrule
\end{tabular}%
}
\end{table}

\section{Conclusion}
This study has demonstrated that a data-centric strategy, specifically the augmentation of a training dataset with negative samples, is a highly effective method for enhancing the reliability of object detection models. By retraining a YOLOv8 model with this refined dataset, we increased detection precision to 94.7\%, mitigating the critical issue of false positives. This transformation of the model's behavior, achieved by teaching it what to ignore, was the key to creating a deployable system. The integration of a secondary classification model further enriches the system's output by providing actionable data on defect severity.

\begin{thebibliography}{9}
\bibitem{yolov8_ultralytics}
Jocher, G., et al. (2023). \textit{Ultralytics YOLOv8}. \url{https://github.com/ultralytics/ultralytics}.

\bibitem{NEURIPS2019_9015}
Paszke, A., et al. (2019). PyTorch: An Imperative Style, High-Performance Deep Learning Library. In \textit{Advances in Neural Information Processing Systems 32}.
\end{thebibliography}

\end{document}
